%=============================================================================!
% 0.2  HOW THE BOOK IS ORGANISED                                              !
%=============================================================================!
\section{How the book is organised}
\label{sec:or-map}

This working edition contains Chapters~0--18 and the trajectory gallery
in Appendix~F. The six-part map in \cref{fig:or-map} also shows the intended
continuation through neural networks and learned propagators, so that the
purpose of the earlier material is visible from the outset.

The six parts follow the order introduced above, and two links between
them shape the reading. The Hamiltonian formulation in Part~I supplies the
conservation laws that TrajCast imposes on its predictions and the geometric
property that keeps the standard integrator stable, so Part~II can locate
exactly the limit on the step that learned propagators aim to lift. Part~IV
ends with a study of lithium in graphite run with a foundation potential,
and its trajectories provide data for the planned comparisons in Part~VI,
where the learned models will be judged on a system the reader has already
simulated.

\begin{figure}[!hb]
\centering
%=============================================================================!
% FIGURE 0.1  CHAPTER MAP                                                     !
% Rows are parts; Parts I, V and VI take two rows each. Solid outline: on     !
% the fast track; dashed: off it. Chapters that the TrajCast chapter draws    !
% on directly carry an accent dot.                                            !
%=============================================================================!
\begin{tikzpicture}[
   x=1cm, y=1cm, font=\small,
   ch/.style={draw=black!40, line width=0.4pt, rounded corners=2pt,
      minimum width=2.7cm, minimum height=0.95cm, align=center,
      text width=2.55cm, inner sep=2pt, fill=white,
      execute at begin node={\hyphenpenalty=10000\exhyphenpenalty=10000}},
   off/.style={ch, dashed},
   goal/.style={ch, draw=accent, line width=0.8pt},
   flow/.style={-{Stealth[length=4.5pt]}, draw=black!40, line width=0.4pt},
   part/.style={anchor=east, align=right, inner sep=0pt},
]
% \partlabel{numeral}{name}: teal lining numeral over a grey small-caps name
\newcommand{\partlabel}[2]{%
   {\color{accent}\liningfigs\normalsize #1}\\[1pt]%
   {\color{black!55}\footnotesize\textsc{#2}}}
% % column centres; rows close within a part and open between parts
\def\ca{2.4}\def\cb{5.6}\def\cc{8.8}\def\cd{12.0}
\def\ra{0}\def\rb{-1.35}\def\rc{-3.0}\def\rd{-4.65}
\def\re{-6.3}\def\rf{-7.95}\def\rg{-9.3}\def\rh{-10.95}\def\ri{-12.3}
\def\lx{0.75}

% PART I
\node[part] at (\lx,{(\ra+\rb)/2}) {\partlabel{I}{mechanics}};
\node[ch]  (c1)  at (\ca,\ra) {1\enskip Motion};
\node[ch]  (c2)  at (\cb,\ra) {2\enskip Newton's laws};
\node[ch]  (c3)  at (\cc,\ra) {3\enskip Work and energy};
\node[ch]  (c4)  at (\cd,\ra) {4\enskip Oscillations};
\node[ch]  (c5)  at (\ca,\rb) {5\enskip Rotation};
\node[ch]  (c6)  at (\cb,\rb) {6\enskip Lagrangian mechanics};
\node[ch]  (c7)  at (\cc,\rb) {7\enskip Hamiltonian mechanics};
\node[ch]  (c8)  at (\cd,\rb) {8\enskip Potential energy surfaces};

% PART II
\node[part] at (\lx,\rc) {\partlabel{II}{integration}};
\node[ch]  (c9)  at (\ca,\rc) {9\enskip Integrators};
\node[ch]  (c10) at (\cb,\rc) {10\enskip Building an MD code};
\node[off] (c11) at (\cc,\rc) {11\enskip Constraints and RESPA};

% PART III
\node[part] at (\lx,\rd) {\partlabel{III}{ensembles}};
\node[ch]  (c12) at (\ca,\rd) {12\enskip Ensembles};
\node[ch]  (c13) at (\cb,\rd) {13\enskip Thermostats};
\node[off] (c14) at (\cc,\rd) {14\enskip Pressure and barostats};

% PART IV
\node[part] at (\lx,\re) {\partlabel{IV}{practice}};
\node[ch]  (c15) at (\ca,\re) {15\enskip Convergence and error};
\node[ch]  (c16) at (\cb,\re) {16\enskip Observables};
\node[ch]  (c17) at (\cc,\re) {17\enskip Lithium in graphite};

% PART V
\node[part] at (\lx,{(\rf+\rg)/2}) {\partlabel{V}{learning}};
\node[ch]  (c18) at (\ca,\rf) {18\enskip Learning as optimisation};
\node[ch]  (c19) at (\cb,\rf) {19\enskip Neural networks};
\node[ch]  (c20) at (\cc,\rf) {20\enskip Symmetry and descriptors};
\node[ch]  (c21) at (\cd,\rf) {21\enskip Spherical harmonics};
\node[ch]  (c22) at (\ca,\rg) {22\enskip Message passing, MACE};
\node[off] (c23) at (\cb,\rg) {23\enskip Training potentials};
\node[off] (c24) at (\cc,\rg) {24\enskip Electrodes and electrolytes};

% PART VI
\node[part] at (\lx,{(\rh+\ri)/2}) {\partlabel{VI}{propagators}};
\node[ch]  (c25) at (\ca,\rh) {25\enskip Time coarsening};
\node[ch]  (c26) at (\cb,\rh) {26\enskip Small propagators};
\node[off] (c27) at (\cc,\rh) {27\enskip Attention and operators};
\node[goal](c28) at (\ca,\ri) {28\enskip TrajCast};
\node[goal,dashed] (c29) at (\cb,\ri) {29\enskip ATOM};
\node[off] (c30) at (\cc,\ri) {30\enskip Outlook for electrodes};

% reading order within rows
\foreach \a/\b in {c1/c2,c2/c3,c3/c4,c5/c6,c6/c7,c7/c8,c9/c10,c10/c11,
                   c12/c13,c13/c14,c15/c16,c16/c17,c18/c19,c19/c20,c20/c21,
                   c22/c23,c23/c24,c25/c26,c26/c27,c28/c29,c29/c30}
   \draw[flow] (\a) -- (\b);
% row wraps
\draw[flow] (c4.south) |- ++(0,-0.2) -| (c5.north);
\draw[flow] (c21.south) |- ++(0,-0.2) -| (c22.north);
\draw[flow] (c27.south) |- ++(0,-0.2) -| (c28.north);

% chapters the TrajCast chapter draws on directly
\foreach \n in {c5,c9,c10,c13,c16,c17,c21,c22,c25,c26}
   \fill[accent] (\n.north east) ++(-0.13,-0.13) circle (1.5pt);
\end{tikzpicture}

\caption{The course map, including the planned Chapters~19--30, in reading
order. An accent dot marks a
chapter on which the TrajCast chapter draws directly; a dashed outline marks
a chapter off the fast track, the shortest path to reading the TrajCast code
with understanding.}
\label{fig:or-map}
\end{figure}
